% فصل سوم: شبکه‌های عصبی آگاه از فیزیک (روش پژوهش)
\chapter{شبکه‌های عصبی آگاه از فیزیک}
\label{ch:method}

در این فصل، چارچوب شبکه‌های عصبی آگاه از فیزیک که روش اصلی این پژوهش است،
به‌تفصیل تشریح می‌شود. ابتدا صورت کلی مسئله و ایده محوری روش بیان می‌گردد
و سپس دو گونه مدل زمان‌پیوسته و زمان‌گسسته برای مسائل مستقیم معرفی می‌شوند.
پس از آن، فرمول‌بندی مسئله معکوس و کشف پارامترها توضیح داده می‌شود و در
پایان، نکات عملی مهم در پیاده‌سازی و آموزش این شبکه‌ها مرور می‌گردد. مبنای
این فصل، مقاله رئیسی و همکارانش \cite{raissi2019} است.

\section{صورت کلی مسئله و ایده محوری}
\label{sec:formulation}

معادلات دیفرانسیل جزئی پارامتری و غیرخطی به شکل کلی زیر را در نظر بگیرید:
\begin{equation}
u_t + \mathcal{N}[u; \lambda] = 0, \quad \mathbf{x} \in \Omega, \quad t \in [0,T],
\label{eq:general-pde}
\end{equation}
که در آن $u(t,\mathbf{x})$ جواب مجهول، $\mathcal{N}[\cdot;\lambda]$ یک
عملگر دیفرانسیل غیرخطی با پارامترهای $\lambda$ و $\Omega$ دامنه مکانی
مسئله است. این صورت کلی، طیف وسیعی از مسائل از قوانین بقا و فرایندهای پخش
تا سامانه‌های واکنش ـ انتشار را در بر می‌گیرد. برای نمونه، معادله برگرز
\eqref{eq:burgers-intro} متناظر با حالت
$\mathcal{N}[u;\lambda] = \lambda_1 u u_x - \lambda_2 u_{xx}$ با
$\lambda = (\lambda_1, \lambda_2)$ است.

با فرض در دست بودن تعدادی مشاهده (احتمالاً نویزی) از سیستم، دو مسئله مجزا
مطرح می‌شود. مسئله نخست، مسئله مستقیم است: با فرض معلوم بودن پارامترهای
$\lambda$، درباره حالت پنهان $u(t,\mathbf{x})$ چه می‌توان گفت؟ مسئله دوم،
مسئله معکوس یا کشف معادله است: کدام پارامترهای $\lambda$ داده‌های مشاهده‌شده
را به بهترین شکل توصیف می‌کنند؟

ایده محوری شبکه‌های آگاه از فیزیک آن است که جواب $u(t,\mathbf{x})$ را با یک
شبکه عصبی عمیق $u_\theta(t,\mathbf{x})$ با پارامترهای $\theta$ تقریب بزنیم
و سپس باقیمانده معادله را به‌صورت زیر تعریف کنیم:
\begin{equation}
f(t,\mathbf{x}) := u_t + \mathcal{N}[u],
\label{eq:residual}
\end{equation}
که در آن $u$ همان خروجی شبکه است. نکته کلیدی اینجاست که چون مشتق‌های
$u_t$، $u_x$، $u_{xx}$ و... همگی با مشتق‌گیری خودکار از روی شبکه به دست
می‌آیند، خودِ $f$ نیز یک شبکه عصبی (با همان پارامترهای $\theta$ ولی توابع
فعال‌ساز متفاوت) محسوب می‌شود؛ به این شبکه، شبکه آگاه از فیزیک گفته می‌شود.
اگر شبکه $u$ به‌خوبی آموزش ببیند، باید در همه نقاط دامنه $f \approx 0$
باشد، یعنی معادله دیفرانسیل ارضا شود. بنابراین کافی است تابع هزینه‌ای
طراحی کنیم که هم فاصله از داده‌ها و هم باقیمانده معادله را هم‌زمان کمینه
کند. نمایی کلی از این چارچوب در شکل \ref{fig:pinn} نشان داده شده است.

\begin{figure}[t]
\centering
\begin{latin}
\begin{tikzpicture}[>=Stealth, node distance=2.2cm,
  block/.style={rectangle, draw, thick, minimum width=2.6cm, minimum height=1.1cm, align=center, font=\small},
  smallblock/.style={rectangle, draw, thick, minimum width=2.2cm, minimum height=1.1cm, align=center, font=\small}]
  \node[smallblock] (input) {Input\\$(t, \mathbf{x})$};
  \node[block, right=1.2cm of input] (net) {Deep neural\\network $u_\theta$};
  \node[smallblock, right=1.2cm of net] (ad) {Automatic\\differentiation};
  \node[block, right=1.2cm of ad] (res) {PDE residual\\$f_\theta = u_t + \mathcal{N}[u]$};
  \node[block, below=1.4cm of ad] (loss) {Loss\\$\mathrm{MSE}_u + \mathrm{MSE}_f$};
  \node[smallblock, left=1.2cm of loss] (opt) {Optimizer\\Adam / L-BFGS};
  \draw[->, thick] (input) -- (net);
  \draw[->, thick] (net) -- (ad);
  \draw[->, thick] (ad) -- (res);
  \draw[->, thick] (net) -- (loss) node[midway, right, font=\small] {$u$ data};
  \draw[->, thick] (res) -- (loss);
  \draw[->, thick] (loss) -- (opt);
  \draw[->, thick, dashed] (opt) -- (net) node[midway, left, font=\small] {update $\theta$};
\end{tikzpicture}
\end{latin}
\caption{نمای کلی چارچوب شبکه عصبی آگاه از فیزیک: شبکه جواب، مشتق‌گیری خودکار، شبکه باقیمانده و حلقه آموزش}
\label{fig:pinn}
\end{figure}

\section{مدل‌های زمان‌پیوسته برای مسئله مستقیم}
\label{sec:continuous}

در مدل زمان‌پیوسته، شبکه ورودی $(t,\mathbf{x})$ را می‌گیرد و مستقیماً
$u(t,\mathbf{x})$ را پیش‌بینی می‌کند؛ به همین دلیل، این مدل‌ها خانواده
جدیدی از تقریب‌زننده‌های تابعی زمانی ـ مکانی به شمار می‌روند. پارامترهای
مشترک شبکه‌های $u$ و $f$ با کمینه‌سازی تابع هزینه زیر آموخته می‌شوند:
\begin{equation}
\mathrm{MSE} = \mathrm{MSE}_u + \mathrm{MSE}_f,
\label{eq:loss-forward}
\end{equation}
که در آن:
\begin{equation}
\mathrm{MSE}_u = \frac{1}{N_u}\sum_{i=1}^{N_u} \left|u(t^i_u,\mathbf{x}^i_u) - u^i\right|^2,
\qquad
\mathrm{MSE}_f = \frac{1}{N_f}\sum_{i=1}^{N_f} \left|f(t^i_f,\mathbf{x}^i_f)\right|^2.
\label{eq:loss-terms}
\end{equation}
در اینجا $\{t^i_u,\mathbf{x}^i_u,u^i\}_{i=1}^{N_u}$ داده‌های آموزشی روی شرط
اولیه و مرز هستند و $\{t^i_f,\mathbf{x}^i_f\}_{i=1}^{N_f}$ نقاط هم‌مکانی\footnote{Collocation Points}
نامیده می‌شوند؛ یعنی نقاطی از دامنه که در آن‌ها صرفاً ارضای معادله (صفر
بودن $f$) اجبار می‌شود، بدون آن‌که مقدار جواب در آن نقاط معلوم باشد.

نکته قابل‌توجه، کارایی داده‌ای این روش است: در مثال‌های مقاله مرجع، تنها با
چند ده یا چند صد نقطه داده اولیه و مرزی، جواب کل دامنه با دقت بالا بازسازی
می‌شود، زیرا بخش $\mathrm{MSE}_f$ مانند یک منظم‌ساز فیزیکی عمل می‌کند و
فضای جواب‌های قابل‌قبول را به‌شدت محدود می‌سازد. در مسئله نمونه این
پایان‌نامه (معادله برگرز)، شبکه $f(t,x)$ شکل صریح زیر را دارد:
\begin{equation}
f := u_t + u u_x - \nu u_{xx},
\label{eq:burgers-f}
\end{equation}
و تابع هزینه دقیقاً همان \eqref{eq:loss-forward} است.

یکی از محدودیت‌های مدل زمان‌پیوسته آن است که وقتی تعداد نقاط هم‌مکانی لازم
زیاد می‌شود (مثلاً در مسائل با دینامیک تند یا ابعاد بالا)، هزینه محاسباتی
آموزش افزایش می‌یابد. برای غلبه بر این مشکل، گونه دوم مدل‌ها معرفی شده
است که در بخش بعد می‌آید.

\section{مدل‌های زمان‌گسسته}
\label{sec:discrete}

در مدل زمان‌گسسته، به‌جای آن‌که شبکه کل بازه زمانی را یکجا بیاموزد، از یک
طرح گام‌زنی زمانی ضمنی از خانواده رانگ ـ کوتا\footnote{Runge-Kutta} با $q$
مرحله استفاده می‌شود. با اعمال صورت کلی روش‌های رانگ ـ کوتا به معادله
\eqref{eq:general-pde} داریم:
\begin{equation}
\begin{aligned}
u^{n+c_i} &= u^n - \Delta t \sum_{j=1}^{q} a_{ij}\, \mathcal{N}[u^{n+c_j}], \quad i=1,\dots,q, \\
u^{n+1} &= u^n - \Delta t \sum_{j=1}^{q} b_j\, \mathcal{N}[u^{n+c_j}],
\end{aligned}
\label{eq:rk}
\end{equation}
که در آن $\{a_{ij}, b_j, c_j\}$ ضرایب طرح رانگ ـ کوتا هستند. در اینجا به‌جای
یک شبکه، $q+1$ خروجی شبکه در نظر گرفته می‌شود که مراحل میانی
$u^{n+c_1},\dots,u^{n+c_q}$ و جواب گام بعدی $u^{n+1}$ را پیش‌بینی می‌کنند و
تابع هزینه، سازگاری این خروجی‌ها با روابط \eqref{eq:rk} و داده‌های دو
لحظه زمانی $t^n$ و $t^{n+1}$ را اجبار می‌کند.

مزیت اصلی این رهیافت آن است که چون طرح رانگ ـ کوتای ضمنی ذاتاً پایدار است،
می‌توان گام‌های زمانی بسیار بزرگ برداشت و حتی با تعداد مراحل زیاد (مثلاً
$q=500$ در برخی مثال‌های مقاله) به دقت زمانی بالایی رسید؛ کاری که در
روش‌های کلاسیک به‌سبب هزینه حل دستگاه‌های غیرخطی بزرگ عملاً ناممکن است.
به بیان دیگر، شبکه عصبی نقش حل‌کننده دستگاه غیرخطی ناشی از گسسته‌سازی
ضمنی را بازی می‌کند. شرح کامل نظریه و انواع طرح‌های رانگ ـ کوتا را می‌توان
در منابع استاندارد روش‌های عددی معادلات دیفرانسیل معمولی یافت
\cite{butcher2016}. در این پایان‌نامه تمرکز پیاده‌سازی بر مدل زمان‌پیوسته
است، اما آشنایی با مدل زمان‌گسسته برای درک کامل چارچوب ضروری بود.

\section{مسئله معکوس و کشف معادله}
\label{sec:inverse}

در مسئله معکوس، علاوه بر جواب $u$، پارامترهای $\lambda$ عملگر دیفرانسیل نیز
مجهول‌اند و باید از روی داده‌ها فراگرفته شوند. خوشبختانه در چارچوب آگاه از
فیزیک، این کار به‌طور طبیعی انجام می‌شود: کافی است $\lambda$ را نیز جزء
پارامترهای قابل‌آموزش شبکه به حساب آوریم و همان تابع هزینه
\eqref{eq:loss-forward} را نسبت به $(\theta, \lambda)$ کمینه کنیم. در اینجا
داده‌های آموزشی، مشاهده‌های پراکنده جواب در سراسر دامنه زمانی ـ مکانی
هستند (نه فقط روی شرط اولیه و مرز) و نقاط هم‌مکانی نیز می‌توانند همان نقاط
داده باشند.

برای مثال، در مسئله کشف معادله برگرز، باقیمانده به شکل زیر در نظر گرفته
می‌شود:
\begin{equation}
f := u_t + \lambda_1 u u_x - \lambda_2 u_{xx},
\label{eq:burgers-f-inv}
\end{equation}
که $\lambda_1$ و $\lambda_2$ از مقدار اولیه دلخواه (مثلاً صفر) شروع می‌کنند
و در جریان آموزش به مقادیر واقعی خود ($\lambda_1=1$ و
$\lambda_2=0.01/\pi$) همگرا می‌شوند. نکته جالب این است که شبکه در این حالت
هم‌زمان دو کار انجام می‌دهد: هم جواب را در کل دامنه بازسازی می‌کند و هم
پارامترهای معادله حاکم را شناسایی می‌نماید. مقاله مرجع نشان می‌دهد که این
شناسایی حتی در حضور نویز در داده‌ها نیز با دقت قابل‌قبولی انجام می‌شود؛
موضوعی که در فصل چهارم، در بخش عملی این پایان‌نامه نیز آزموده خواهد شد.

تفاوت مهم مسئله معکوس با روش‌هایی مانند \lr{SINDy} (که در فصل دوم معرفی
شد) در این است که در اینجا نیازی به کتابخانه‌ای از جملات کاندید نیست و
مشتق‌های لازم نیز به‌جای تفاضل محدودِ روی داده‌های نویزی، با مشتق‌گیری
خودکار از روی شبکه‌ای هموار محاسبه می‌شوند؛ همین امر روش را در برابر نویز
مقاوم‌تر می‌سازد.

\section{نکات عملی در پیاده‌سازی و آموزش}
\label{sec:practical}

تجربه مقاله مرجع و پژوهش‌های بعدی نشان می‌دهد که موفقیت آموزش شبکه‌های آگاه
از فیزیک به چند انتخاب مهم بستگی دارد که در ادامه خلاصه می‌شوند.

\textbf{معماری شبکه.} در بیشتر مثال‌ها از پرسپترون‌های چندلایه نسبتاً کوچک
(مثلاً ۴ تا ۸ لایه پنهان با ۲۰ تا ۵۰ نورون در هر لایه) با فعال‌ساز تانژانت
هیپربولیک استفاده می‌شود و معمولاً نیازی به منظم‌سازهای اضافی مانند افت
تصادفی\footnote{Dropout} نیست، زیرا خودِ جمله فیزیکی تابع هزینه نقش
منظم‌ساز را ایفا می‌کند.

\textbf{بهینه‌ساز.} راهبرد استاندارد، شروع آموزش با چند هزار گام آدام و سپس
ادامه آن با \lr{L-BFGS} تا همگرایی است. آدام در ابتدا سریع پیش می‌رود و از
نقاط زینی عبور می‌کند، در حالی که \lr{L-BFGS} در نزدیکی بهینه، دقت نهایی را
به‌طور چشمگیری بهبود می‌بخشد. در این پایان‌نامه نیز دقیقاً همین راهبرد دو
مرحله‌ای به کار رفته است.

\textbf{نقاط هم‌مکانی.} تعداد و توزیع نقاط هم‌مکانی اثر مستقیمی بر دقت و
هزینه آموزش دارد. معمولاً چند هزار نقطه که به‌صورت تصادفی یکنواخت (یا با
نمونه‌برداری لاتین هایپرمکعب\footnote{Latin Hypercube Sampling}) از دامنه
انتخاب می‌شوند کفایت می‌کند. در مسائلی که جواب در ناحیه خاصی تغییرات تندی
دارد، تمرکز بیشتر نقاط در آن ناحیه می‌تواند مفید باشد.

\textbf{مقیاس‌بندی.} از آنجا که ورودی‌های شبکه (زمان و مکان) و خروجی آن
ممکن است در بازه‌های متفاوتی باشند، نرمال‌سازی ورودی‌ها به بازه‌ای مانند
$[-1,1]$ معمولاً به همگرایی سریع‌تر و پایدارتر کمک می‌کند. در پیاده‌سازی
این پایان‌نامه، دامنه مسئله چنان انتخاب شده که نیازی به مقیاس‌بندی اضافی
نباشد.

در جدول \ref{tab:models} دو گونه مدل معرفی‌شده در این فصل از نظر ویژگی‌های
اصلی مقایسه شده‌اند.

\begin{table}[t]
\caption{مقایسه مدل‌های زمان‌پیوسته و زمان‌گسسته در چارچوب آگاه از فیزیک}
\label{tab:models}
\centering
\singlespacing
\begin{tabular}{|c|c|c|}
\hline
\textbf{ویژگی} & \textbf{مدل زمان‌پیوسته} & \textbf{مدل زمان‌گسسته} \\
\hline
ورودی شبکه & $(t, \mathbf{x})$ & $\mathbf{x}$ در یک لحظه زمانی \\
\hline
خروجی شبکه & $u(t,\mathbf{x})$ & مراحل رانگ ـ کوتا و $u^{n+1}$ \\
\hline
داده آموزشی & شرط اولیه و مرزی & جواب در دو لحظه $t^n$ و $t^{n+1}$ \\
\hline
مزیت اصلی & سادگی و یکپارچگی & گام زمانی بزرگ و پایدار \\
\hline
کاربرد در این پایان‌نامه & پیاده‌سازی شده & فقط بررسی نظری \\
\hline
\end{tabular}
\end{table}

\section{جمع‌بندی فصل}
\label{sec:ch3-summary}

در این فصل، چارچوب شبکه‌های عصبی آگاه از فیزیک تشریح شد: ایده تعریف شبکه
باقیمانده با مشتق‌گیری خودکار، تابع هزینه ترکیبی داده و فیزیک، مدل‌های
زمان‌پیوسته و زمان‌گسسته برای مسئله مستقیم، و فرمول‌بندی مسئله معکوس برای
کشف پارامترها. در فصل بعد، مدل زمان‌پیوسته برای معادله برگرز پیاده‌سازی و
نتایج عددی آن گزارش و تحلیل می‌شود.
